\scanpage{130}

% L1–2
$\lambda\mapsto\lambda^2$ sur le centre $\mathbb G_m$). Ici \deleted{pour} $n$ impair on trouve encore

% L3, slanting left margin
% slanted
\margin{commencer par ceci avant de prendre $O(q)$}

% L3–4
\uline{Prop} $O'(q)\overset{\sim}{\longleftarrow}\underbrace{SO(q)}_{\substack{\parallel\,(\mathrm{def}^{\mathrm n})\\O'(q)\cap\mathrm{Sl}(E)}}\times\mathbb G_m$.

% L5
qui équivaut au ---

% L6–11
\uline{Cor} (Pour $\varepsilon=\pm1$ fixé) La catégorie des modules quasi-orthogonaux \inserted{$(E,L,q)$ de rg $n$ impair $n$} sur $S$ est équiv.lente à la cat. des \deleted{\unsure{cples}} syst. $(E_0,q_0,\omega_0,\Delta)$, où $(E_0,q_0,\omega_0)$ est un module \inserted{spéc.} orth pour $\varepsilon$, de rang $n$, $\Delta$ est un module inv. À un système on associe

% L7–11, sideways left margin
% margin note written sideways
\margin{Dégager Thm plus \unsure{rigoureux|rigureux}}

% L12
\[
 E=E_0\otimes\Delta,\quad\inserted{L=\Delta^{\otimes 2},}\quad
 \varepsilon q_0^{\deleted{\unsure{\otimes\Delta|\otimes\Delta^{\otimes2}}}}:E_0\otimes\Delta\longrightarrow\Delta^{\otimes 2}=L,
\]

% L12, right margin
\margin{$\left(\begin{aligned}
 &\text{donc }\det E\simeq\overbrace{\det E_0}^{\mathrm{bse}}\otimes\Delta^n\\
 &\simeq\deleted{\det E_0\otimes{}}\Delta^{2\nu}\otimes\Delta\\
 &\simeq\deleted{\det E_0\otimes{}}L^\nu\otimes\Delta
\end{aligned}\right)$}

% L13
Le foncteur quasi-inv. associe à $(E,L,q)$

% L14, left margin
\margin{$n=2\nu+1$}

% L14
\[
 \Delta=\inserted{L^{\otimes-\nu}}\otimes\det E,\quad
 E_0=E\otimes\Delta^{-1},\quad q_0\deleted{\unsure{=|:}}\ \varepsilon q\otimes\Delta^{-1}:E_0=E\otimes\Delta^{-1}\longrightarrow L\Delta^{-2}
\]

% L15
\[
 \text{or }L\Delta^{-2}\simeq L\otimes L^{2\nu}\otimes(\det E)^{-2}\simeq L^n\otimes(\det E)^{-2}\simeq\mathcal O_S
 \quad\text{par }(\unsure{1.3|4.3})\quad\deleted{\text{par }(1.3)\quad\mathcal O_S}
\]

% L16
$\omega_0\in\Gamma(\det E_0)$ déduit de

% L17–20
\[
\begin{aligned}
 \det E_0&\simeq\underbrace{\det E}_{\Delta}\otimes\Delta^{-n}
 \deleted{=\Delta\,\unsure{\Delta^n\Delta^{2\nu+1}}\simeq L^\nu\Delta\quad\text{or}}\\
 &\deleted{\simeq\Delta L^\nu\Delta^{-n}}\\
 &\deleted{\det E_0\simeq\det E\otimes\Delta E^{\unsure{-\nu}}}\\
 &\simeq\det E\,L^{n\nu}(\det E)^{-n}=L^{n\nu}(\det E)^{-2\nu}\\
 &=(L(\det E)^{-2})^\nu=\mathcal O_S^{\otimes\nu}\simeq\mathcal O_S\qquad\unsure{\text{cf Thm}|\text{etc.}}
\end{aligned}
\]

% lower left
% archive sheet number: 129
